\begin{abstract}
Data quality rules in recurring batch pipelines must accommodate legitimate change without accepting persistent errors as normal. Fixed thresholds can raise false alarms under distribution drift, whereas unrestricted adaptation can weaken error detection. This article presents and evaluates an adaptive and cost-aware monitoring framework that combines numerical threshold adaptation, gated drift estimation, historical admission of large language model (LLM) rules, guarded category learning, and risk-based check selection. Separate drift and failure signals assign each batch one of four labels: stable, drift, degradation, or drift with degradation. Evaluation uses three public datasets, seventeen controlled scenarios, and ten paired test seeds under an internally frozen protocol, supplemented by two error-labelled benchmarks and a sixteen-year chronological replay. Four-class macro-F1 ranges from 0.947 to 0.966, although the same gate with fixed thresholds performs comparably on two datasets. Adaptation reduces drift-related false positives on flight data but loses up to 29 percentage points of recall under persistent contamination; its false-positive criterion is not supported on retail data. Fixed statistical rules combined with admitted LLM rules achieve 99.6--100\% recall on injected domain-impossible values. A separate per-value category check admits all twelve valid novel values and rejects eleven of twelve invalid values. Risk-based selection saves 2--11\% of checks, but provides little total runtime benefit after profile statistics are reused. The findings identify where adaptation, domain knowledge, and selective execution contribute to monitoring quality, and the conditions that limit those contributions.
\end{abstract}
